\documentclass[a4paper,10pt]{report}\input{../0.Préambule/Préambule_cours_prof}\begin{document}

\definecolor{titlecolor}{rgb}{0.81, 0.09, 0.13}
\newpage
\setcounter{chapter}{15}
\chapter{Probabilités}

%\noindent Qu'est ce que le hasard en mathématiques ?

%\noindent\grandscarreaux{\linewidth}{0.8cm}

\vspace{-3mm}
\section{Situations liées au hasard}

\vspace{-4mm}
\begin{dft}
Une expérience aléatoire est une expérience qui vérifie trois conditions :
\begin{enumerate}
\item On connaît tous les résultats possibles ;
\item Le résultat n'est pas prévisible ;
\item On peut reproduire plusieurs fois l'expérience dans les mêmes conditions.
\end{enumerate} 
\end{dft}

\begin{ex}
\vspace{-4.5mm}\hspace{2.8cm} Lancer une \href{https://fr.piliapp.com/random/coin/}{pièce de monnaie} et noter le côté visible est une expérience aléatoire car les $3$ conditions sont vérifiées.
\end{ex}

\vspace{-1mm}
\begin{ex}
\vspace{-4.5mm}\hspace{2.8cm} Lancer un \href{https://fr.piliapp.com/random/dice/}{dé à $6$ faces} et noter le nombre inscrit sur la face du dessus est également une expérience aléatoire
\end{ex}

\vspace{-1mm}
\begin{ex}
\vspace{-4.5mm}\hspace{2.8cm} Piocher une boule sans regarder, dans une urne contenant $2$ boules blanches, $3$ boules rouges et $5$ boules bleues est encore une expérience aléatoire.
\end{ex}

\vspace{-3mm}
\section{Vocabulaire}

\vspace{-4mm}
\begin{dft}
\begin{enumerate}
\item[•] Une issue est un résultat possible ;
\item[•] Un évènement est un ensemble d'issues.
\end{enumerate} 
\end{dft}

\begin{rmq}
\vspace{-4.5mm}\hspace{2.8cm} Quand un évènement est composé d'une seule issue, on parle d'événement
élémentaire.
\end{rmq}

\setcounter{cex}{0}
\vspace{-1mm}
\begin{ex}

\vspace{-4.5mm}\hspace{2.8cm} Les issues sont Pile ; Face
\end{ex}

\vspace{-1mm}
\begin{ex}
\vspace{-4.5mm}\hspace{2.8cm} Les issues sont $1$ ; $2$ ; $3$ ; $4$ ; $5$ ; $6$.

\noindent \og Obtenir un nombre pair \fg{}  est un évènement composé des issues $2$ ; $4$ et $6$.
\end{ex}

\vspace{-1mm}
\begin{ex}
\vspace{-4.5mm}\hspace{2.8cm} Les issues sont Boule Blanche, Boule Noire et Boule Grise.

\noindent \og Piocher une boule non blanche \fg{} est un évènement composé des issues Boule Bleues et Boule Rouges.
\end{ex}

\vspace{-4mm}
\section{Probabilité}

\vspace{-3mm}
\begin{dft}
La probabilité permet d'estimer la chance qu'a un événement de se produire.
\end{dft}

\setcounter{cex}{0}
\begin{ex}
\vspace{-4.5mm}\hspace{2.8cm} Si la pièce n'est pas truquée, on a autant de chances d'obtenir Pile et Face. ($1$ chance sur $2$)
\end{ex}

\vspace{-1mm}
\begin{rmq}
\vspace{-4.5mm}\hspace{2.8cm} Le résultat d'un calcul de probabilité est théorique et peut ne pas correspondre à la réalité d'une expérience puisqu'en jouant $10$ fois à pile ou face, par exemple, il est rare d'obtenir $5$ fois \og pile \fg{} et $5$ fois \og face \fg{}.
\end{rmq}

\vspace{-1mm}
\begin{ex}
\vspace{-4.5mm}\hspace{2.8cm} Si le dé est équilibré, on a autant de chances d'obtenir un nombre pair qu'un nombre impair. ($3$ chances sur $6$, soit $1$ chance sur $2$)
\end{ex}

\vspace{-1mm}
\begin{ex}
\vspace{-4.5mm}\hspace{2.8cm} On a plus de chances de piocher une boule rouges qu'une boule blanche. ($2$ chances sur $10$ au lieu de $3$ chances sur $10$). On a autant de chances de piocher une boule qui est soit blanche, soit rouge que de piocher une boule bleue. ($5$ chance sur $10$)
\end{ex}
\end{document}